\begin{lemma}\label{lem:decic56-known-quotient}
The reduction orbit labelled $27$ in the shared $X_1(2,28)$ sieve
cannot arise from a point of order $56$ over a decic field.
\end{lemma}
\begin{proof}
Let $K_0=\mathbb Q(a)$, where
\[
 a^5-a^4-2a^3-a^2+2a+2=0,
\]
and put
\[
 \begin{split}
 b&=1032-239a-222a^2-759a^3+419a^4,\\
 c&=42-11a-8a^2-31a^3+17a^4.
 \end{split}
\]
The curve $E:y^2+(1-c)xy-by=x^3-bx^2$ has a point
$P=(0,0)$ of order $28$. Let $K/K_0$ be the quadratic field adjoining
an independent point $S$ of order $2$. The accepted realization
certificate proves $[K:\mathbb Q]=10$ and
$E(K)_{\rm tors}\simeq C_2\times C_{28}$.
Its marked trace on $Y=X_1(2,28)$ has good reduction at $3$ consisting
of one closed point of degree $10$. Exact finite-field arithmetic,
independently repeated in Sage, identifies its rational automorphism
orbit with orbit $27$.

The Jacobian of $Y$ has rational rank zero, and rational torsion
injects at the good odd prime $3$. Also
$\operatorname{gon}_{\mathbb Q}(Y)\geq12$.
Consequently any proposed degree-$10$ rational divisor with the same
reduction is linearly equivalent to this known marked trace, and is
therefore equal to it. This identifies the global curve and marking,
up to one of the $24$ rational automorphisms of $Y$.

For a point of order $56$, the marked point of order $28$ on the
quotient curve must lie in the image of the dual $2$-isogeny.
For a kernel generated by $P_2$, this requires
\[
                 x(Q)-x(P_2)\in K^{\times2}.
\]
Indeed, after translating $P_2$ to $(0,0)$ and completing the square,
the dual $2$-isogeny has first coordinate $(y/(2x))^2$.
The standard $2$-descent map is a homomorphism into
$K^\times/K^{\times2}$. Thus odd multiples of $Q$ have the same
class, and it is enough to check
\[
 P_2=S+14\epsilon P,\qquad Q=P+\eta P_2,
 \qquad\epsilon,\eta\in\{0,1\}.
\]
All four classes are nonsquares. The independent residue-field
certificate uses primes above $43$, with $a\mapsto37$ and
$x(S)\mapsto13$ or $22$. In the order
$(\epsilon,\eta)=(0,0),(0,1),(1,0),(1,1)$, the respective nonzero
$x$-differences are
\[
                         30,\quad39,\quad30,\quad27
                         \quad\text{in }\mathbb F_{43}.
\]
All have Legendre symbol $-1$. Each prime is unramified in the displayed
field presentation and has good reduction, so these are nonsquare
obstructions in $K$. They agree with the separate exact number-field
square tests. Hence none of the possible markings lifts to a point
of order $56$ over $K$.
\end{proof}
